\section{A concrete route toward quasi-polynomial recognition}\label{sec:research}

The locality theorem separates a search-theoretic cost that is now controlled
from geometric obligations that are not. This section states those obligations
as questions with testable intermediate targets. A change in a parameter or
source class should be reflected in the theorem statement, rather than hidden
inside an empirical success rate.

\subsection{The sufficient coverage theorem we would need}

Here is one precise form of a future global result. It allows a controlled
list of sources, rather than assuming that a single height seed works on
every diagram.

\begin{proposition}[Conditional recognition consequence]\label{prop:conditional}
Suppose there is an algorithm that, from an $n$-crossing knot diagram,
constructs at most $n^{O(\log n)}$ authenticated exterior triangulations with
specified integral cocycles and region budgets. Assume its construction
cost is $n^{O(\log n)}$, all source sizes and marking bit lengths are
polynomial in $n$, every region budget satisfies
$c=O(\log n)$ and $r+U=O(\log^2 n)$, and a complete component-disc predicate
has polynomial bit complexity on each resulting endpoint. Finally assume
the following uniform coverage property: every unknot diagram has a
compressing-disc endpoint in at least one of the resulting region families.
Then exhaustive region search gives an $n^{O(\log n)}$ unknot-recognition
algorithm.
\end{proposition}
\begin{proof}
Theorem~\ref{thm:local} bounds each search by $n^{O(\log n)}$; multiplying
by the number of sources preserves this bound. Source authentication and
component-disc replay make every positive verdict sound. The uniform
coverage property guarantees a positive endpoint on every unknot. Therefore,
after all the finite searches finish without one, the diagram is knotted.
The stated predicate cost ensures that this exhaustion is an algorithm,
not a search with unresolved endpoint queries.
\end{proof}

The proposition identifies exactly how a positive coverage theorem would
also justify negative decisions by exhaustion. No separate short negative
certificate is required for that logical consequence. The current artifact
does not establish uniform coverage and therefore does not take that final
negative step. A few short successful traces, or a coverage theorem only for
an already recognized subfamily of unknots, would be insufficient.

An alternative program can work through a certified hierarchy. In that case
one must bound the number, size and marking cost of all subsequent search
sources, as in \eqref{eq:outer-sum}, and prove the outer hierarchy complete.
A local simplification that can restart without a global decreasing measure
does not yet supply that bound.

\subsection{Priority 1: geometric locality and an honest obstruction search}

\paragraph{1. Which source families admit a logarithmic or squared-logarithmic region?}
Fix a concrete diagram-to-exterior construction, its shelling policy and
its cocycle normalization. Ask whether every solid torus in this source
family admits a disc endpoint with $r+U=O(\log^2 t)$, first with one connected
allowed region and then with $O(\log t)$ components. The theorem should
specify whether the cocycle is fixed, optimized by vertex potentials, or
selected from a bounded list. Changing that choice changes the coverage
problem.

A useful initial milestone is a proof for one explicit infinite family of
unknot diagrams that survives the existing shortcut stages and creates
nontrivial interior move choices. Merely testing a family whose initial
coherent seed is already a disc would establish little about geometric
search. The cancelled-braid empty-region miss and the eight-tetrahedron
obstruction are starting controls with different failure mechanisms.

\paragraph{2. Can the required region or upward budget be proved large?}
For increasing triangulated solid tori, compute the minimum allowed-region
size and minimum total upward count separately. Certified lower bounds
require complete exhaustion of the smaller family, with no capped endpoint
queries. Such computations can refute an overly strong conjecture on a
finite example, while a family construction could rule out a proposed
uniform logarithmic theorem. Useful adversarial sources include long layered
regions, separated interacting defects, and relabelled copies that expose
heuristics based on arbitrary row order.

The target is not to optimize a single scalar $r+U$ before understanding
both terms. A small permitted spatial region may still need many expansion
events, and a trace with few upward events can consume a large initial
region. The current count makes their different geometric meanings visible.

\paragraph{3. Can hierarchy structure produce the needed regions?}
Relate the initial consumption footprint to the product regions,
multi-surfaces and controlled decompositions in the hierarchy program
\cite{LackenbyTalk,Lackenby2026}. A useful lemma would take one precisely
specified reduction object, produce a legal marked trace, and bound its
consumed original cells, number of components and total expansions. The
proof must include the representation of the surface and its marking after
cutting. An existence theorem for a topological simplification without this
translation leaves the algorithmic parameters unbounded.

The successful reduction object may be much smaller than the normal surface
it encodes. Binary multiplicity is therefore compatible with a small region
theorem; enumerating all normal sheets to locate that region is not.
Conversely, the height-growth lemma alone controls arithmetic and gives
no spatial localization. These two estimates should be proved separately
and combined only at the final cost accounting.

\paragraph{4. Can small excess height be converted to a small total budget?}
Pachner-graph experiments often constrain maximum excess tetrahedron count
\cite{BurtonPachner}, whereas the theorem here constrains total upward
events. A conversion theorem would need to remove redundant excursions or
bound essential returns. Commuting disjoint moves removes order redundancy
but does not eliminate a sequence of dependent expansion--collapse cycles.
An exact state-based loop-removal argument may help with shortest traces;
its state count and marking dependence must be included. A bound on upward
bursts has the same issue and should not be substituted for $U$.

\subsection{Priority 2: stronger complete geometric search}

\paragraph{5. Can separators replace total active-region size?}
The commitment bound charges every freely assigned incarnation. In a long
trace whose interactions have a small separator, a dynamic program might
reuse completed subregions instead. A promising parameter is the width of
the dependency graph of events, with edges whenever one event consumes a
cell created or constrained by another. The first theorem needed is a
bounded interface description: which boundary face maps, local cocycle
values, surviving identities and surface data suffice to compose partial
traces? Without that description, invoking treewidth only renames the
state-explosion problem.

A concrete target is a family with $r=\Theta(t)$ but logarithmic interface
width, together with a finite compatibility relation and a proved state
count polynomial in the marking bit length. This would extend the useful
class beyond small initial footprints.

\paragraph{6. Which additional local relations safely reduce trace classes?}
The present sleep relation uses disjoint supports. Burton's overlapping
composite relations suggest larger reductions \cite{BurtonPachner}. Each
candidate needs a marked relative-boundary proof and an event-identity rule
that remains stable under other reductions. Test completeness against the
naive oracle before evaluating speed. A raw triangulation transposition table
must include remaining resources and continuation restrictions, or supply a
dominance proof. Equal unmarked endpoints alone do not justify dropping
one search state.

Region enumeration itself offers a modest immediate target: share exact
endpoint-predicate results across allowed regions while retaining separate
continuations. The proof should distinguish the terminal surface property,
which can be reused, from the permissions to consume original cells,
which differ. This avoids repeatedly checking the source endpoint and other
common states without changing the searched union.

\paragraph{7. Can the framework include the other Pachner moves?}
For unmarked formal star replacements, a tentative alphabet would have
six edge ports, four face ports, four vertex ports, one whole-cell port
and idle, for sixteen symbols. If $u_{23},u_{14}$ count expansions and
$d_{32},d_{41}$ count contractions, set $E=u_{23}+3u_{14}$. The tetrahedron
balance gives
\[
 d_{32}+3d_{41}\le t+E.
\]
The birth count then obeys
\[
 t+3u_{23}+4u_{14}+2d_{32}+d_{41}\le3t+5E.
\]
This accounting suggests an analogous constant-alphabet theorem, but the
formal-star legality, commutation and marking rules must be proved before
asserting it for an implementation. In particular, a $1$--$4$ move introduces
an interior vertex whose cocycle height needs a canonical choice or an
explicitly bounded search. The current artifact does not implement that
extension.

\paragraph{8. Can positive traces be made smaller without weakening replay?}
The explicit-snapshot certificate is polynomial for one bounded trace, but
repeats unchanged surrounding tetrahedra at every move. A persistent face-map
representation or a shared-prefix certificate DAG could remove that
duplication. The independent verifier should accept a stream of local changes,
check the same relative-boundary facts, and bound its memory. Storing all
endpoints of an exponential search is a different problem: a compact proof
of one success does not compress the whole search tree.

\subsection{Priority 3: exact LP discovery and normal-surface structure}

\paragraph{9. Can a genuine triangulation realize a superconstant screening gain?}
The algebraic separation has unit coefficients and sparse equations, but
normal matching matrices also encode very specific face and disc-type
incidences. Construct a family of actual compact triangulations for which
the ordered baseline makes asymptotically more deletion calls than local
witness screening. A weaker useful first target is a rigorous lower bound
on the number of redundant baseline tests in a layered or glued family,
with a matching positive point and independently checked topology. The
finite solid-torus timings do not establish such an asymptotic family.

\paragraph{10. What geometric hypotheses bound the support antichain?}
Proposition~\ref{prop:exponential-antichain} rules out a small-antichain
argument from sparsity and unit coefficients alone. Possible additional
hypotheses include a controlled decomposition of the triangulation, a
bounded number of quad types interacting with the boundary, or a restricted
rank-one cone. The desired theorem should bound the number of useful
minimal supports or provide a small subcollection covering all deletion
queries encountered by a specified branch policy. Bounding all supports
may be stronger than the algorithm needs.

\paragraph{11. Can first-node LPs be accelerated with exact dual lifting?}
The capped figure-eight controls make this a direct implementation priority.
An equality-row basis can reduce a redundant matrix before simplex. A
certificate of the row transformation allows a dual of the reduced problem
to be lifted to the original equations. Fraction-free elimination and sparse
factorization are natural candidates, but their exact coefficient growth,
per-mask rank changes and reconstruction cost need measurements and proofs.

Such an optimization can preserve mathematical soundness while changing
the exact serialized multiplier and hence falls outside the byte-preserving
mode proved here. It should have its own compatibility contract. More
ambitiously, a polynomial-bit-complexity LP method can replace the pivot
engine while retaining exact independently checked primal or dual outputs;
one must count its recovery and verification work, not merely invoke a
real-arithmetic LP bound. Warm starts are another practical direction,
but need careful basis validation after coordinates are deleted.

\paragraph{12. How should positive and negative evidence interact?}
Positive supports safely answer any further zero restriction containing
them. A negative dual has the opposite issue when columns are restored:
its inequality must be checked on every newly allowed column. A two-sided
cache could reuse both forms of evidence with full matrix checks. Its
branch policy and final proof may change, so compare it against the present
byte-preserving producer rather than silently weakening that guarantee.
Useful quantities include exact validation cost, cache hit location, and
the number of current-node LPs avoided.

\paragraph{13. Can propagation backdoors be bounded after geometric preprocessing?}
The incoming dual-certificate work isolates quadrilateral propagation
backdoors. The locality search may expose triangulations with smaller
backdoors even when it does not directly find a disc. A joint theorem
would bound both the cost of reaching such a triangulation and the number
of unresolved choices afterwards. The endpoint predicate would then be
algebraic complexity reduction rather than immediate disc discovery. It
must remain source-bound, and the outer choice of endpoints must not
reintroduce an uncontrolled number of LP searches.

\subsection{Priority 4: evaluation that can disprove the proposed story}

\paragraph{14. What does a matched source-level cohort show?}
Create a frozen cohort stratified by original diagram crossings, exterior
tetrahedron count, known shortcut success and initial coherent-disc status.
Every arm should receive identical face pairings and marking choices.
Compare time to a verified positive or negative result, the proportion
remaining inconclusive under explicit caps, total upward moves, region size,
endpoint-query time, LP pivots and certificate size. Independent relabellings
and matched constructor versions are essential controls for order-sensitive
heuristics. The present fifteen LP cases and six geometric control sizes
are useful mechanism tests, not such a representative recognition study.

\paragraph{15. Which empirical claims survive fixed-resource reproduction?}
Repeat paired timing in an isolated environment, report all censored runs,
and retain the raw per-pair observations. Separate source construction,
search, terminal verification and later independent replay when their cost
matters. A speedup that disappears after including required verification is
not an end-to-end gain. A future benchmark should also include the unfavorable
small cases and node-LP bottlenecks already observed here.

The highest-value next result is a geometric theorem bounding the parameters
in Proposition~\ref{prop:conditional}, or a rigorous obstruction showing
that a proposed bound fails. The next implementation result is exact
current-node LP acceleration with a checked reconstruction map. Both have
clear acceptance criteria and address limitations that the present proofs
and experiments identify directly.
